% ============================================================
%  B201 Computer Science Lab -- Portfolio Project Report
%  Author : Parvin Riyah | June 2026
%  Font   : Times New Roman (mathptmx), 12 pt body
%  Colour : Deep-purple / violet scheme
% ============================================================

\documentclass[12pt, a4paper]{report}

% ---- Fonts & encoding --------------------------------------
% Latin Modern (lmodern) is a modern Type1 replacement for Computer Modern.
% NOTE: For Times New Roman on a full install, replace the three lines below with:
%   \usepackage[T1]{fontenc}  \usepackage[utf8]{inputenc}  \usepackage{mathptmx}
\usepackage[T1]{fontenc}
\usepackage[utf8]{inputenc}
\usepackage{lmodern}

% ---- Page geometry -----------------------------------------
\usepackage[
  a4paper,
  top=2.5cm, bottom=2.5cm,
  left=3cm,  right=2.5cm
]{geometry}

% ---- Colours -----------------------------------------------
\usepackage{color}
\definecolor{deepPurple}{rgb}{0.176, 0.106, 0.412}
\definecolor{midPurple}{rgb}{0.357, 0.176, 0.557}
\definecolor{lightPurple}{rgb}{0.608, 0.349, 0.816}
\definecolor{paleLavender}{rgb}{0.929, 0.906, 0.965}
\definecolor{accentViolet}{rgb}{0.486, 0.227, 0.929}
\definecolor{darkText}{rgb}{0.102, 0.102, 0.180}
\definecolor{white}{rgb}{1,1,1}

% ---- Coloured table rows -----------------------------------
\usepackage{colortbl}

% ---- Tables ------------------------------------------------
\usepackage{array}
\usepackage{tabularx}

% ---- Spacing & paragraph -----------------------------------
\linespread{1.25}
\setlength{\parskip}{7pt}
\setlength{\parindent}{0pt}

% ---- Header / footer (applied after titlepage) -------------
\usepackage{fancyhdr}

% ---- Verbatim for code blocks ------------------------------
\usepackage{verbatim}

% ---- Hyperlinks (works with lmodern + T1) ------------------
\usepackage{hyperref}
\hypersetup{
  colorlinks=true,
  urlcolor=midPurple,
  linkcolor=darkText,
  pdfborder={0 0 0}
}

% ---- Chapter / section title formatting -------------------
%  We redefine the internal LaTeX macros to colour headings.
\makeatletter

% Chapter heading -- H1: 16 pt bold deep-purple
\def\@makechapterhead#1{%
  \vspace*{18pt}%
  {\parindent\z@ \raggedright \normalfont
   \fontsize{16}{20}\bfseries\color{deepPurple}%
   \thechapter.\hspace{0.5em}#1
   \par\nobreak
   \vskip 10pt
   \textcolor{midPurple}{\rule{\linewidth}{0.8pt}}%
   \vskip 8pt
  }%
}

% Unnumbered chapter heading (used by \chapter*)
\def\@makeschapterhead#1{%
  \vspace*{18pt}%
  {\parindent\z@ \raggedright \normalfont
   \fontsize{16}{20}\bfseries\color{deepPurple}%
   #1\par\nobreak
   \vskip 10pt
   \textcolor{midPurple}{\rule{\linewidth}{0.8pt}}%
   \vskip 8pt
  }%
}

% Section -- H2: 14 pt bold mid-purple
\renewcommand\section{%
  \@startsection{section}{1}{\z@}%
    {-14pt plus -4pt minus -2pt}%
    {6pt}%
    {\normalfont\fontsize{14}{18}\bfseries\color{midPurple}}%
}

% Subsection -- H3: 12 pt bold accent-violet
\renewcommand\subsection{%
  \@startsection{subsection}{2}{\z@}%
    {-10pt plus -3pt minus -1pt}%
    {4pt}%
    {\normalfont\fontsize{12}{16}\bfseries\color{accentViolet}}%
}

\makeatother

% ============================================================
\begin{document}

% ============================================================
%  PAGE 1 -- COVER / ASSESSMENT SUBMISSION FORM
% ============================================================
\begin{titlepage}
\thispagestyle{empty}

% ---- Top purple band (simulated with coloured rule) --------
\noindent
\colorbox{deepPurple}{%
  \parbox[t][2.2cm]{\textwidth}{%
    \vspace{0.5cm}
    \centering
    \textcolor{white}{\large\textbf{Gisma University of Applied Sciences}}\\[4pt]
    \textcolor{white}{\normalsize B201 Computer Science Lab}
  }%
}

\vspace{-0.3cm}
\noindent\textcolor{midPurple}{\rule{\linewidth}{3pt}}\\[-2pt]
\noindent\textcolor{lightPurple}{\rule{\linewidth}{1.5pt}}

\vspace{1.2cm}

% ---- Report title block ------------------------------------
\begin{center}
  {\fontsize{22}{28}\selectfont\bfseries\color{deepPurple}
    Portfolio Project Report}\\[14pt]
  \textcolor{midPurple}{\rule{0.55\textwidth}{1pt}}\\[20pt]
  \textcolor{lightPurple}{\rule{0.35\textwidth}{0.6pt}}
\end{center}

\vspace{1.1cm}

% ---- Assessment Submission Form ----------------------------
\begin{center}
  {\fontsize{13}{17}\selectfont\bfseries\color{deepPurple}
    Assessment Submission Form}
\end{center}

\vspace{0.4cm}

\renewcommand{\arraystretch}{1.6}
\noindent
\begin{tabularx}{\textwidth}{%
  >{\bfseries\color{darkText}\raggedright\arraybackslash}p{5.8cm}%
  >{\raggedright\arraybackslash}X}
\rowcolor{paleLavender}%
Student Number & Parvin Riyah \hspace{0.4em}--\hspace{0.4em} GH1047690 \\
\multicolumn{2}{>{\small\itshape\color{midPurple}\arraybackslash}l}{%
  \quad (If this is group work, please include the student numbers of all group participants)}\\[2pt]
\rowcolor{paleLavender}%
Assessment Title    & Portfolio Project Report \\
Module Code         & B201D \\
\rowcolor{paleLavender}%
Module Title        & Computer Science Lab \\
Module Tutor        & William Baker Morrison \\
\rowcolor{paleLavender}%
Date Submitted      & 26.06.2026 \\
\end{tabularx}

\vspace{0.9cm}

% ---- Declaration of Authorship -----------------------------
\begin{center}
  {\fontsize{13}{17}\selectfont\bfseries\color{deepPurple}
    Declaration of Authorship}
\end{center}

\vspace{0.35cm}

\begin{center}
\begin{minipage}{0.88\textwidth}
\fontsize{12}{17}\selectfont\color{darkText}
\centering

I declare that all material in this assessment is my own work except where
there is clear acknowledgement and appropriate reference to the work of
others.

\vspace{10pt}

I fully understand that the unacknowledged inclusion of another person's
writings or ideas or works in this work may be considered plagiarism and
that, should a formal investigation process confirm the allegation, I would
be subject to the penalties associated with plagiarism, as per Gisma
Business School, University of Applied Sciences' regulations for academic
misconduct.

\vspace{20pt}

Signed.\enspace Parvin Riyah\enspace
Date.\enspace 26.06.2026

\end{minipage}
\end{center}

% ---- Bottom purple band ------------------------------------
\vfill
\noindent\textcolor{midPurple}{\rule{\linewidth}{1.5pt}}\\[-3pt]
\noindent\colorbox{deepPurple}{%
  \parbox[t][1.0cm]{\textwidth}{%
    \vspace{0.3cm}\centering
    \textcolor{white}{\small Portfolio Project Report \textbar{} Parvin Riyah \textbar{} June 2026}%
  }%
}

\end{titlepage}

% ============================================================
%  PAGE 2 -- TABLE OF CONTENTS
% ============================================================
\newpage
\pagenumbering{roman}
\setcounter{page}{1}
\pagestyle{fancy}
\fancyhf{}
\renewcommand{\headrulewidth}{0.6pt}
\fancyhead[L]{\small\textit{\textcolor{midPurple}{B201 Computer Science Lab}}}
\fancyhead[R]{\small\textit{\textcolor{midPurple}{Portfolio Project Report}}}
\fancyfoot[C]{\small\thepage}

\tableofcontents
\clearpage

% ============================================================
%  MAIN BODY
% ============================================================
\pagenumbering{arabic}
\setcounter{page}{1}
\renewcommand{\chaptername}{}

% ============================================================
%  1. INTRODUCTION
% ============================================================
\chapter{Introduction}

This report details the creation, structural design, and technical workflow
of my personal computer science portfolio, developed as a core component of
the B201 Computer Science Lab module at Gisma University of Applied Sciences.
The goal of this portfolio is to build a cohesive, recognisable professional
identity that showcases my software engineering skills and academic progress
to prospective employers.

The portfolio ecosystem is built around three pillars: a custom-designed
personal website, a structured GitHub repository hosting the code, and a
clean, high-fidelity CV typeset in LaTeX. Together, these assets provide
a transparent look into my technical proficiencies and development practices.

\section*{Links}

\begin{itemize}
  \item \textbf{LinkedIn:}\enspace
        \href{https://www.linkedin.com/in/parvin-riyah-13724a39a}%
             {https://www.linkedin.com/in/parvin-riyah-13724a39a}
  \item \textbf{GitHub:}\enspace
        \href{https://github.com/ParvinRiyah757}%
             {https://github.com/ParvinRiyah757}
    \item \textbf{GitHub Repostary:}\enspace
        \href{https://github.com/ParvinRiyah757/Parvin-Riyah_Portfolio.git}%
             {https://github.com/ParvinRiyah757/Parvin-Riyah_Portfolio.git}
  \item \textbf{Website:}\enspace
        \href{https://github.com/ParvinRiyah757/Parvin-Riyah_Portfolio.git}%
             {https://github.com/ParvinRiyah757/Parvin-Riyah_Portfolio.git}
  \item \textbf{CV:}\enspace
        \href{https://github.com/ParvinRiyah757/Parvin-Riyah_Portfolio.git}%
             {https://github.com/ParvinRiyah757/Parvin-Riyah_Portfolio.git}
  \item \textbf{Report:}\enspace
        \href{https://github.com/ParvinRiyah757/Parvin-Riyah_Portfolio.git}%
             {https://github.com/ParvinRiyah757/Parvin-Riyah_Portfolio.git}
\end{itemize}

% ============================================================
%  2. PORTFOLIO WEBSITE STRUCTURE
% ============================================================
\chapter{Portfolio Website Structure}

The portfolio website is engineered as a seamless single-page experience,
built from scratch using clean HTML and CSS. It is divided into the
following intuitive sections:

\begin{itemize}
  \item \textbf{Hero\,/\,Header:} Introduces my name and profile as a
        computer science student, immediately establishing the site's
        personal branding.

  \item \textbf{About Me:} A brief overview detailing my background, core
        values like pragmatism and perfectionism, and my current path.

  \item \textbf{Technical Skills:} A visually striking layout showcasing my
        proficiencies in Python, HTML/CSS, and SQL/MySQL.

  \item \textbf{Language Metrics:} A dedicated section displaying my language
        proficiencies, highlighted by my C1 English level, using custom
        geometric indicators.

  \item \textbf{Project Cards:} Deep dives into my development projects,
        complete with architecture summaries and direct links to their source
        code repositories.

  \item \textbf{Education:} A timeline tracking my academic steps, including
        secondary education and my current Bachelor of Computer Science degree
        at Gisma University of Applied Sciences.

  \item \textbf{Footer:} Clean, accessible links for direct contact and
        professional networking.
\end{itemize}

The page also features an easily accessible download button inside the
header/about section, allowing visitors to instantly grab the PDF version of
my LaTeX CV.

% ============================================================
%  3. GITHUB REPOSITORY STRUCTURE
% ============================================================
\chapter{GitHub Repository Structure}

To maintain an organized development workspace, the portfolio repository is
structured cleanly, separating core frontend code from underlying project
assets:

\noindent\textcolor{midPurple}{\rule{\linewidth}{0.5pt}}
{\small
\begin{verbatim}
parvin-riyah-portfolio/
|-- Report        # Compiled LaTeX project report
|-- Cv        # Compiled LaTeX CV for direct download
+-- Website    #Main HTML5 structure and layout rules
    
\end{verbatim}
}
\noindent\textcolor{midPurple}{\rule{\linewidth}{0.5pt}}

While smaller files are stored directly within the primary portfolio
repository, larger database components---such as my E-sports Database
Management System---are hosted in their own dedicated GitHub repositories to
preserve distinct version control histories and separate documentation.

% ============================================================
%  4. DESIGN DECISIONS
% ============================================================
\chapter{Design Decisions}

\section{Website Design \& The Blue Aesthetic}

I chose to build this website using pure HTML5 and CSS3 without touching
heavy frameworks or external libraries. I wanted to prove that I understand
core web technologies from the ground up, keeping the site lightning-fast
and highly maintainable.

For the visuals, I wanted the site to feel deeply authentic to who I am.
Blue is my favourite colour, so I used that as the foundation for my
personal branding. I designed a modern, ``cosmic'' dark mode using a deep,
immersive blue background (\verb|#0B0D1F|) paired with vibrant accent tones.
To make it highly functional, I built a smooth fluid theme toggle using CSS
custom properties, allowing users to shift into a clean, crisp off-white
light mode layout (\verb|#F7F7FC|) effortlessly.

\section{Responsive Layout Strategy}

Unlike rigid, desktop-only layouts, I prioritised a fully responsive design.
By combining CSS Flexbox and Grid with fluid typography math
(\verb|clamp()|), the entire layout scales seamlessly. Whether a hiring
manager reviews my work on a wide desktop monitor or a smartphone on the go,
the content shifts naturally without breaking boundaries or creating awkward
horizontal scrolls.

\section{CV Design via LaTeX}

For my CV, I bypassed standard word processors and wrote it from scratch
in LaTeX. LaTeX gives me absolute control over typesetting, spacing,
and font weight hierarchies. By manually structuring the document layout
instead of relying on a crowded pre-made template, I ensured that the file
remains perfectly parsable by automated Applicant Tracking Systems (ATS)
while staying highly readable and visually satisfying for human recruiters.

% ============================================================
%  5. TOOLS AND TECHNOLOGIES
% ============================================================
\chapter{Tools and Technologies}

\begin{itemize}
  \item \textbf{HTML5 \& CSS:} Used to construct the entire presentation
        layer from scratch, utilising CSS variables, grids, and flexboxes to
        establish a fluid UI.

  \item \textbf{LaTeX:} Selected for drafting a clean, publication-grade
        r\'{e}sum\'{e}, giving me precise control over typography and margins.

  \item \textbf{Python:} The core language used to write my Hangman game,
        handling the underlying terminal game loops, data processing, and
        state evaluation.

  \item \textbf{MySQL:} Used to architect my E-sports management database
        and the backend state storage for the Hangman game, focusing on
        relational integrity and query optimisation.

  \item \textbf{Git \& GitHub:} Used for local version control and remote
        code hosting, tracking development iterations cleanly through clear
        commit histories.
\end{itemize}

% ============================================================
%  6. STRENGTHS AND WEAKNESSES
% ============================================================
\chapter{Strengths and Weaknesses}

\section{Strengths}

\begin{itemize}
  \item \textbf{Fluid Responsiveness:} The website adapts instantly to any
        screen size, preserving layout integrity from mobile devices to
        desktop monitors.

  \item \textbf{Strong Personal Branding:} The blue-driven cosmic aesthetic
        stands out immediately, creating a memorable visual theme that
        reflects my personal style.

  \item \textbf{Clean Code Architecture:} Writing the site without heavy
        frameworks keeps the codebase lightweight, highly readable, and
        exceptionally fast to load.

  \item \textbf{ATS-Optimised CV:} The LaTeX r\'{e}sum\'{e} strikes a
        great balance between clean professional design and machine-readable
        data structures.
\end{itemize}

\section{Weaknesses}

\begin{itemize}
  \item \textbf{Early-Stage Project Volume:} As an ongoing student, my
        portfolio currently highlights two main projects, which limits the
        total breadth of technologies shown.

  \item \textbf{Terminal-Based Gameplay:} The Hangman game executes directly
        inside the terminal, which demonstrates solid logical skills but lacks
        immediate visual feedback for non-technical users.
\end{itemize}

% ============================================================
%  7. FUTURE IMPROVEMENTS
% ============================================================
\chapter{Future Improvements}

\begin{itemize}
  \item \textbf{Advanced UI Transitions:} I plan to add subtle CSS
        micro-animations to cards and menu selections to make navigation feel
        even more dynamic and responsive.

  \item \textbf{Expanding the Project Suite:} As I move forward in my
        computer science degree, I plan to integrate more complex projects,
        specifically focusing on full-stack web applications and interactive
        algorithms.

  \item \textbf{Upgrading to a Graphical Interface:} I would love to
        refactor the Python Hangman game by replacing the terminal display
        with a custom GUI (using libraries like Pygame or Tkinter) to give it
        more visual impact.

  \item \textbf{Integrated Contact Form:} Replacing the basic email links
        with a built-in contact form would allow recruiters to message me
        directly without leaving the website.
\end{itemize}

\end{document}
